\documentclass[UTF8,10pt,a4paper,fontset=fandol]{ctexart}
\usepackage[margin=17mm,headheight=15pt]{geometry}
\usepackage{amsmath,amssymb,booktabs,longtable,array,tabularx,enumitem,fancyhdr,xcolor,tikz}
\usetikzlibrary{arrows.meta,positioning,trees}
\usepackage[unicode,colorlinks=true,linkcolor=blue!45!black,urlcolor=blue!45!black]{hyperref}
\hypersetup{pdftitle={CSC5010 期中开卷参考资料},pdfauthor={个人复习资料}}
\setlength{\parindent}{0pt}\setlength{\parskip}{4pt}
\setlist{nosep,leftmargin=1.6em}
\renewcommand{\arraystretch}{1.18}
\pagestyle{fancy}\fancyhf{}\fancyhead[L]{CSC5010 · 期中开卷参考}\fancyhead[R]{\nouppercase{\leftmark}}\fancyfoot[C]{\thepage}
\renewcommand{\sectionmark}[1]{\markboth{#1}{}}
\newcommand{\src}[1]{\par{\footnotesize\color{black!65}来源：#1。所有页码为 PDF 文件页码。}\par}
\newcommand{\ex}[1]{\subsection{#1}}
\newcommand{\checkit}[1]{\par\textbf{检查：}#1\par}
\newcommand{\method}[1]{\par\textbf{何时用：}#1\par}
\newcommand{\ans}[1]{\boxed{#1}}
\newcommand{\blank}{\square}
\newcolumntype{Y}{>{\raggedright\arraybackslash}X}
\begin{document}
\begin{center}
{\LARGE\bfseries CSC5010 期中开卷参考资料}\\[5pt]
{\large 定义 · 公式 · 手算流程 · 作业与课件例题}\\[4pt]
2026 年 11 月 17 日期中考试复习版\quad 编写日期：2026-11-13
\end{center}
\textbf{使用方式：}先按下表找题型，再翻目录所示页码。各章保留计算条件、中间表和分数结果。\textbf{“教学变式”是编写者为演示方法自拟的题，不是教师原题。}资料供个人复习与纸质开卷使用，遵守课件不得向班外传播的限制。

\textbf{范围：}课程开始至 10 月 29 日的 Search、Logic Representation、Probabilistic Reasoning；HW1--HW5、HW6 P1--P5。\textbf{Bayesian Network、HW6 P6--P7 不收录。}考试通知为纸质开卷，禁止计算器，计算可保留分数。本册不把背景扩展当作教师押题。

\par\noindent\begin{tabularx}{\textwidth}{p{.32\textwidth}Y r}
\toprule 想查什么 & 关键词 / 入口 & 页码\\\midrule
Agent、理性、XOR & Rational action / state modeling & \pageref{sec:intro}\\
队列、BFS、DFS、时间戳 & BFS / DFS / parent / backtrack & \pageref{sec:search}\\
拓扑序、迭代加深 & DAG / finish time / IDDFS & \pageref{ex:topo}\\
带权最短路、松弛 & Dijkstra / relaxation & \pageref{ex:dijkstra}\\
g、h、f、改父节点、重开 & A* / admissible / consistent / reopen & \pageref{sec:astar}\\
前瞻、树节点、Minimax & branching / value backup & \pageref{sec:games}\\
剪枝、开放线评分 & alpha--beta / evaluation & \pageref{ex:ab}\\
MRV、degree、LCV、删域 & CSP / forward checking & \pageref{sec:csp}\\
弧一致、局部修复 & AC-3 / min-conflicts & \pageref{ex:ac}\\
必要充分、真值表、范式 & logic / NNF / CNF / DNF & \pageref{sec:logic}\\
归结、满足赋值、规则链 & resolution / SAT / Horn & \pageref{ex:resolution}\\
组合、二项、多项、平滑 & binomial / multinomial / Laplace & \pageref{sec:prob}\\
后验、全概率、分类 & Bayes / Naive Bayes & \pageref{sec:bayes}\\
联合边缘、独立检验 & joint / marginal / chi-square & \pageref{sec:markov}\\
转移、矩阵、稳态、等待 & Markov / stationary / first passage & \pageref{ex:transition}\\
题号索引与易错条件 & HW1--6 / 来源映射 & \pageref{sec:crossref}\\\bottomrule
\end{tabularx}\par

\tableofcontents
\clearpage
\section{概念与状态空间}\label{sec:intro}
\src{课件 01 p.13--15、27--34；02 p.4--28；HW1 INTRO1--3}
\par\noindent\begin{tabularx}{\textwidth}{p{.22\textwidth}Y}
\toprule 概念 & 可直接用于短答的解释\\\midrule
理性行动 Rational action & 在已知信息、可用行动与资源条件下，选择预期最能达成目标的行动。依据行动时的信息评价，随机结果不保证每次成功；理性不等于自私或道德好坏。\\
Agent & 接收环境输入，按目标作决策，产生行动或输出的系统；可以是机器人，也可以是软件。\\
四种 AI 定义 & 像人思考、理性思考、像人行动、理性行动。本课偏重理性行动。\\
强/弱 AI & 一般智能的目标 / 解决特定任务的系统。特定任务成功不自动说明具有一般理解。\\
中文屋 & 符号操作产生看似正确的行为，是否等同理解；属于行为与理解关系的哲学讨论。\\
人不擅长的任务 & 例：持续、精确检查百万条记录的重复与规则冲突。说明注意力/记忆负担，以及算法如何用明确规则处理。\\
XOR & 00,11 为 0；01,10 为 1，线性不可分。隐藏层可构造 $h_1=x_1\lor x_2,h_2=x_1\land x_2$，输出 $h_1\land\neg h_2$；纯线性层复合仍线性。\\\bottomrule
\end{tabularx}\par

\ex{搜索建模与计数模板}
\method{题目给现实情境，要求定义搜索或估算规模。}
依次写：状态 $s$、起点 $s_0$、合法动作、后继函数、目标条件、一步代价和路径代价。状态应包含决定未来动作所需信息。
\[
\text{m 个变量，每个 k 值：}\ |S|=k^m;\qquad
\text{固定 b 叉树：第 d 层 }b^d,\quad
N_{0:d}=\frac{b^{d+1}-1}{b-1}.
\]
$b=1$ 时总数 $d+1$；有早停、去重、不等分支时应逐层计数。

\textbf{密码锁例：}3 位、每位 0--9，共 $10^3=1000$ 状态；每步只改变一位 $\pm1$，最多 6 个邻居。每位循环时，从 $x$ 到 $y$ 的最少操作数为
\[
\sum_{i=1}^3\min(|x_i-y_i|,10-|x_i-y_i|).
\]
不循环时为 $\sum_i|x_i-y_i|$。三位十值是三维，不是十维。

\textbf{8-puzzle：}九元组含空格 0；空格和上下左右数字交换，通常每步代价 1。角/边中/中心合法动作数 2/3/4；排列上界 $9!$，固定可解起点可达 $9!/2$。搜索树按历史计节点，原状态图按排列合并。
\checkit{若两个路径到同一棋盘，它们是两个搜索树节点、一个不同状态。BFS 的发现顺序、扩展次数、最短路径、frontier 大小不是同一个量。}

\clearpage\section{BFS、DFS、拓扑序、IDDFS 与 Dijkstra}\label{sec:search}
\src{课件 03 p.3--52；HW2 BFS1、DFS1--5、Dijkstra1}
\par\noindent\begin{tabularx}{\textwidth}{lYYY}
\toprule 方法 & 取点规则 & 常见最优条件 & 手算记录\\\midrule
BFS & FIFO 队列 & 相同步代价 & 队列、层、距离、parent\\
DFS & 栈/递归深入 & 一般不最优 & 递归栈、进入/退出时间\\
IDDFS & 深度限额 0,1,2,… & 相同步代价的最浅解 & 每轮到达/扩展数\\
Dijkstra/UCS & 最小当前 $g$ & 非负边权 & 距离、parent、确定集合\\\bottomrule
\end{tabularx}\par
显式邻接表图 BFS/DFS 时间 $O(V+E)$、空间 $O(V)$。有限图去重时两者可遍历全部可达点；无限深度 DFS 不保证终止。UCS 无限空间常需有限分支且边代价有正下界。堆 Dijkstra 常记 $O((V+E)\log V)$。

\ex{BFS 队列、parent 与最短路径（教学变式）}
无向边 $AB,AC,BD,BE,CE$；起点 A，邻接字母序，入队即标记。
\begin{center}\begin{tabular}{c|l|l|l}\toprule
出队&新发现&本轮后队列&新 parent\\\midrule
A&B,C&B,C&B$\leftarrow$A,C$\leftarrow$A\\
B&D,E&C,D,E&D$\leftarrow$B,E$\leftarrow$B\\
C&无&D,E&不重复发现 E\\D&无&E&--\\E&无&空&--\\\bottomrule
\end{tabular}\end{center}
距离 $d(A)=0,d(B)=d(C)=1,d(D)=d(E)=2$。访问序 A,B,C,D,E；到 E 路径按 parent 反转为 A--B--E。\textbf{同层次序依赖邻接顺序，层数不变。}

\ex{八数码 BFS 计数（HW2 BFS1，p.1）}
初始棋盘 $\left[\begin{smallmatrix}2&8&3\\1&6&4\\7&\blank&5\end{smallmatrix}\right]$。根深度 0；不同状态去重。
\begin{center}\begin{tabular}{c|rrrrrr}\toprule
深度&0&1&2&3&4&5\\\midrule
本层状态&1&3&5&10&14&28\\\bottomrule
\end{tabular}\end{center}
第二层 $2(2-1)+(4-1)=5$，因排除立即回父状态。题设第 5 层查一半：
\[
N=1+3+5+10+14+\frac{28}{2}=\ans{47}.
\]
\checkit{根与目标计入；这是已检查状态数量。若问已生成/扩展节点，需按具体目标检测时机另算。}

\ex{DFS 时间戳与回溯（HW2 DFS1--2）}
时钟每进入或退出顶点加 1。完整 n 顶点遍历产生 2n 个时间。祖先区间包含后代区间：$d(u)<d(v)<f(v)<f(u)$。

\textbf{DFS1 p.2：}按原图课件选择，树边 $uv,vt,ts,tr,rw$，发现 $u,v,t,s,r,w$。
\begin{center}\begin{tabular}{c|rrrrrr}\toprule
顶点&u&v&t&s&r&w\\\midrule
$d$&1&2&3&4&6&7\\$f$&12&11&10&5&9&8\\\bottomrule
\end{tabular}\end{center}
返回 $s\rightsquigarrow t$ 后转入 r，重新分叉点 t；无向非树边 su、wv 连向祖先。无向 DFS 只有树边和回边；有向 DFS 还可能有前向、交叉边。

\textbf{DFS2 p.3：}按右侧原图、字母邻接序，D 是一个顶点，不因左示意树重复标签而重复访问。
\begin{center}\begin{tabular}{c|rrrrrrrr}\toprule
顶点&S&A&B&C&E&D&F&G\\\midrule
$d/f$&1/16&2/15&3/14&4/5&6/13&7/8&9/12&10/11\\\bottomrule
\end{tabular}\end{center}
重新分叉点 B,E；函数返回的每一步也可叫 backtrack，应在答题时声明口径。

\textbf{DFS3 p.4 网格续走：}保留所给前缀，余下邻居按行列标签升序这一约定，可得分叉点 $(4,1),(4,2),(4,4)$。新增路线为 $41\to51\to61\to62$（死路），返回 $42\to43\to44\to34\to24$（死路），再返回 $44\to54$（目标）。题目未固定余下方向顺序，不同声明可产生不同路线。最终路径只保留 parent 链，不包括走入后又退出的死路。

\textbf{DFS4 p.5 有向图：}从 v 开始，余下根/邻接字母序，发现序
\[
v,w,s,q,t,x,z,y,r,u.
\]
根为 v,q,r；退出序 $s,w,v,z,x,y,t,q,u,r$。回边 $s\to v,z\to x,y\to q$；交叉边 $q\to s,q\to w,u\to y,r\to y$。

\ex{拓扑排序（HW2 DFS5 + 教学小图）}\label{ex:topo}
\method{有向无环图中，要求每条边 $u\to v$ 的 u 在 v 之前。}
方法一：完整 DFS 后按 $f$ 递减排序；方法二：反复输出入度 0 点并删出边。存在有向环则没有拓扑序。

教学图 $A\to B,A\to C,B\to D,C\to D$，字母序 DFS：A 1/8，B 2/5，D 3/4，C 6/7；逆退出序 $A,C,B,D$。$A,B,C,D$ 也合法。

HW2 p.6 的一组合法序：
\[
\ans{p,n,o,s,m,r,y,v,x,w,z,u,q,t}.
\]
\checkit{逐条检查 $\operatorname{pos}(u)<\operatorname{pos}(v)$；拓扑序通常不唯一。有向 DFS4 图含环，不可用其逆退出序冒充拓扑序。}

\ex{迭代加深的重复计数（课件 03 p.43；教学变式）}
DLS 到深度 L 截断；cutoff 不等于无解。IDDFS 重跑 L=0,1,…,d，每轮重置该轮记录。
\[
N_{\mathrm{IDDFS}}=\sum_{i=0}^d(d-i+1)b^i.
\]
b=2,d=3，各轮到达节点 $1,3,7,15$，共 $26$；BFS 到深度 3 共 15。若只数真正扩展且不扩展限制层，IDDFS 对应 $0+1+3+7=11$ 次；\textbf{计数先写口径}。树搜索内存约 $O(bd)$。

\ex{Dijkstra 全表（HW2 Dijkstra1，p.7）}\label{ex:dijkstra}
\method{边权非负，要求逐轮松弛与最短路径。}
\[
g(v)\leftarrow\min\{g(v),g(u)+w(u,v)\};\quad
\text{严格改善时同步设置 }parent(v)=u.
\]
每轮取所有未确定点中 $g$ 最小者；目标\textbf{取出确定}才可停止。原题为无向图，特别注意 $w(b,h)=1$。
\begin{center}\small\begin{tabular}{c|rrrrrrrrr}\toprule
确定点&a&b&h&g&f&c&i&e&d\\\midrule
初始&0&$\infty$&$\infty$&$\infty$&$\infty$&$\infty$&$\infty$&$\infty$&$\infty$\\
a&0&4&8&$\infty$&$\infty$&$\infty$&$\infty$&$\infty$&$\infty$\\
b&0&4&5&$\infty$&$\infty$&12&$\infty$&$\infty$&$\infty$\\
h&0&4&5&6&$\infty$&12&12&$\infty$&$\infty$\\
g&0&4&5&6&8&12&12&$\infty$&$\infty$\\
f&0&4&5&6&8&12&12&18&22\\
c&0&4&5&6&8&12&12&18&19\\
i&0&4&5&6&8&12&12&18&19\\
e&0&4&5&6&8&12&12&18&19\\
d&0&4&5&6&8&12&12&18&19\\\bottomrule
\end{tabular}\end{center}
parent：b$\leftarrow$a；h,c$\leftarrow$b；g,i$\leftarrow$h；f$\leftarrow$g；e$\leftarrow$f；d$\leftarrow$c。$g(d):22\to19$ 因 $12+7<8+14$。
\checkit{到 d 路径 a--b--c--d，边和 $4+8+7=19$；到 e 路径 a--b--h--g--f--e，边和 18。12 的平局按字母选 c 再 i。}

\clearpage\section{A* 与启发函数}\label{sec:astar}
\src{课件 04 p.2--3、11、17、20--33；HW2 AStar1--3}
\[
\text{Greedy: }f=h;\qquad\text{Dijkstra: }f=g;\qquad\text{A*: }f=g+h.
\]
$g$ 为当前最好已知累计代价；$h$ 为剩余代价估计；每轮从\textbf{整个 OPEN}取最小 f。后继候选 $g'=g(n)+w(n,m)$，严格更小时更新 g、f、parent 和队列优先级；需要时重开 CLOSED 节点。\textbf{目标出队才停止，不能一生成目标就停。}

\par\noindent\begin{tabularx}{\textwidth}{lY}
\toprule 条件 & 判别与使用\\\midrule
Admissible 可采纳 & $0\le h(n)\le h^*(n)$，$h(G)=0$。不能高估真实最短剩余代价。\\
Consistent 一致 & 每条有向边 $h(u)\le w(u,v)+h(v)$；无向图查两个方向，即 $|h(u)-h(v)|\le w$。\\
一致的作用 & 沿路径 f 非递减；在通常有限非负图上可以不重开已关闭节点。\\
仅可采纳 & 树搜索或支持更优路径 reopen 的图搜索可保最优；永久封闭 CLOSED 一般不够。\\\bottomrule
\end{tabularx}\par

\ex{八数码启发式（课件 04 p.3）}
\[
S=\begin{pmatrix}7&2&4\\5&\blank&6\\8&3&1\end{pmatrix},\qquad
G=\begin{pmatrix}\blank&1&2\\3&4&5\\6&7&8\end{pmatrix}.
\]
错位块数 $h_1=8$；不数空格。Manhattan
\[
h_2=\sum_{k=1}^8(|r_k-r_k^*|+|c_k-c_k^*|)
=3+1+2+2+3+2+2+3=\ans{18}.
\]
每次移动一数字一格，故两者均是下界且一致；以题目给的目标为准。地图直线距离适用于同单位的距离边权，时间权重需相应时间下界。

\ex{原图 AStar1：h(B)=16（HW2 p.8）}
\textbf{足以复算的无向边：}AX17、XB16、XC11、CL20、CW11、WB12、WK10、WT9、BT17、BS7、SH6。启发值如下：
\[
\begin{array}{c|rrrrrrrrrr}n&A&X&B&C&L&W&K&T&S&H\\\hline
h&30&27&16&17&33&9&19&0&17&22\end{array}
\]
记 $n(g,f)$。初始 OPEN=$\{A(0,30)\}$。
\begin{center}\small\par\noindent\begin{tabularx}{\textwidth}{cY Y}\toprule
出队&本轮计算&更新后 OPEN，按 f 排序\\\midrule
A&X: $0+17=17,f=44$&X(17,44)\\
X&B: $17+16=33,f=49$；C: $17+11=28,f=45$&C(28,45),B(33,49)\\
C&W: $28+11=39,f=48$；L: $28+20=48,f=81$&W(39,48),B(33,49),L(48,81)\\
W&T: $39+9=48,f=48$；K: $39+10=49,f=68$；经 W 到 B 为51，拒绝&T(48,48),B(33,49),K(49,68),L(48,81)\\
T&取出目标，停止&--\\\bottomrule\end{tabularx}\par\end{center}
\[
\ans{A\to X\to C\to W\to T},\quad\ans{17+11+11+9=48}.
\]
Greedy 会先 B 再 T，总长 50。\textbf{课件 04 p.11 的局部数字没有累计前缀，手算必须用上表的 g。}

\ex{AStar2：h(B)=6 后更新 OPEN（HW2 p.9）}
\begin{center}\small\par\noindent\begin{tabularx}{\textwidth}{cY Y}\toprule
出队&计算&OPEN\\\midrule
A&产生 X(17,44)&X(17,44)\\
X&B(33,39),C(28,45)&B(33,39),C(28,45)\\
B&W(45,54),T(50,50),S(40,57)&C(28,45),T(50,50),W(45,54),S(40,57)\\
C&W: $45\to39,f=48$；L(48,81)&W(39,48),T(50,50),S(40,57),L(48,81)\\
W&T: $50\to48$；K(49,68)&T(48,48),S(40,57),K(49,68),L(48,81)\\
T&取出目标&停止\\\bottomrule\end{tabularx}\par\end{center}
最优路径/代价仍 48。新 h 可采纳但不一致：$h(X)=27>16+h(B)=22$。本题改善的是 OPEN 中的 W、T，\textbf{本次并未发生 CLOSED 中更优路径重开}。

\ex{reopen 的独立演示（教学变式）}
有向边 S$\to$A=3，S$\to$B=1，B$\to$A=1，A$\to$G=3。$h(S)=h(A)=h(G)=0,h(B)=4$，可采纳但 $4>1+0$ 不一致。
\begin{center}\begin{tabular}{c|l}\toprule
出队&结果\\\midrule
S&A(3,3), B(1,5)\\
A&G(6,6)，A 关闭\\
B&改善 A 为 g=2，必须重开\\
A&G 更新到 g=5\\
G&最优 S--B--A--G，总代价 5\\\bottomrule\end{tabular}\end{center}
永久跳过已关闭 A 会错误返回 6。\checkit{比较改善看 g；选下个节点看 f；返回路径看 parent。}

\ex{多边形障碍（HW2 AStar3，p.10）}
连续位置状态无限；若只取障碍顶点加起终点，共 V+2 个不同节点，连接不穿障碍内部的可见直线段，得 visibility graph。最短路在自由空间的无障碍折点可拉直；必要转折落在障碍顶点。路径由可见线段构成，可能沿部分边界，\textbf{不要求全程沿多边形边}。

\clearpage\section{博弈、Minimax 与剪枝}\label{sec:games}
\src{课件 05 p.16--18、32--67、73--74；HW3 Game1--5}
\textbf{约定：}X=MAX、O=MIN；分数一直按 MAX 视角；ply 为一个玩家行动一次。终局立即停止；X 赢 +1，O 赢 $-1$，和 0。
\[
V(s)=\begin{cases}U(s)&\text{终局}\\E(s)&\text{到截止深度且非终局}\\
\max_{s'}V(s')&\text{轮 MAX}\\\min_{s'}V(s')&\text{轮 MIN}\end{cases}
\]

\ex{分支数、前瞻评估与终局叶数}
固定分支 b、深度 d：只评估底层叶 $b^d$；评估所有非根则 $\sum_{i=1}^db^i$；所有节点含根再加 1。

\textbf{HW3 Game1：}红棋两车 3+2，两马 2+2，两相 2+2，两仕 1+1，帅 1，两炮 13+10，五兵 5，共 $\ans{44}$。右炮包含隔一个炮架的远端吃子。一步叶评估 44；固定 b 的三步叶评估 $44^3=\ans{85184}$，所有非根 $44+1936+85184=87164$。

\textbf{HW3 Game5：}根深度 0，按不合并棋盘/对称的\textbf{合法落子历史}计数。叶只能深度 5--9：5、7 层新终局 X 胜 +1，6、8 层 O 胜 $-1$，9 层 X 胜或和。第 5 层：
\[
\underbrace{8}_{\text{X 获胜线}}\underbrace{3!}_{\text{X 次序}}
\underbrace{\binom62 2!}_{\text{O 位置与次序}}=\ans{1440}.
\]
第 6 层 optional：$8\cdot3!\binom63 3!-12\cdot3!3!=\ans{5328}$；扣去 X 第 5 手已赢的历史。第 9 层 optional 的合法树枚举值为 X 胜 81792、和 46080，共 127872。$9!$ 包含提前终局后仍继续下的非法历史，不是精确合法对局数。

\ex{完整 Minimax 回传与同树剪枝（教学变式）}
根 MAX，三个 MIN 子节点的叶分别 (3,5,6)、(2,9,4)、(4,7,8)。
\[
\min(3,5,6)=3,\quad\min(2,9,4)=2,\quad\min(4,7,8)=4;
\quad V(root)=\ans{4}.
\]
根选第三个分支。左到右剪枝：首分支返回 3，根 $\alpha=3$；第二分支先看 2，$\beta=2\le\alpha$，剪 9,4；第三分支需完整看完，返回 4。共看 7 个叶，根值不变。\textbf{剪枝后的第二分支只需知道上界 2，不能据此声称所有未看叶都等于 2。}

\ex{alpha--beta 的更新与课件两棵树}\label{ex:ab}
初始 $\alpha=-\infty,\beta=+\infty$。MAX 更新 $\alpha=\max(\alpha,v)$；MIN 更新 $\beta=\min(\beta,v)$；$\alpha\ge\beta$ 时剪\textbf{当前节点未访问的孩子}。好排序：MAX 先大、MIN 先小；根值与合法排序无关，剪枝位置和数量有关。

\textbf{课件 05 p.66 上图：}根 MAX，下层 MIN，再 MAX，叶按对分组 $(3,4),(1,2)\mid(7,8),(5,6)$。MAX 层 4,2,8,6；MIN 层 2,6；根 6。左到右\textbf{没有叶子被剪}，需看 8 个叶。

\textbf{同页下图：}$(6,5),(8,7)\mid(2,1),(3,4)$。左 MIN 先得 6；第二 MAX 看到 8 已达 $\beta=6$，剪 7。根 $\alpha=6$；右 MIN 第一 MAX 返回 2，使 $\beta=2\le6$，剪最后整组 3,4。看 6,5,8,2,1，共 5 叶，根 6。

\ex{井字棋开放线评分与两步前瞻（HW3 Game4）}
\[
E(s)=\#\text{无 O 的获胜线}-\#\text{无 X 的获胜线}.
\]
总共八条线；空线同时算两边，已有双方棋子的线两边都不算。\textbf{终局值优先}，可用远大于普通评分的 $\pm M$。

初始 X 中心、O 上中。X 开放线 6，O 开放线 4，E=2。X 先选候选动作，再对每个 O 回应算 E，先取 MIN，再在 X 的候选中取 MAX。
\begin{center}\begin{tabular}{l|l|r}\toprule
X 候选&图中 O 回应分数&MIN\\\midrule
左上角&2,1,3,1,2,1&1\\左中&1,0,2,0,1,1&0\\
左下角&2,2,3,1,2,2&1\\下中&1,1,0（对称代表）&0\\\bottomrule
\end{tabular}\end{center}
选角，根值 1；左右对称得四个角并列。并列后偏选下角属于额外 tie-break，不是严格 Minimax 改变。

\textbf{HW3 Game3：}所给部分树：先角子 MIN 回传 $\min(0,1,1)=0$，先中心 $\min(0,1)=0$，根 $\max(0,0)=0$。这些值针对图示分支；缺省分支不能自动当不存在。\textbf{Game2 续棋：}必须交替、落空格、成三连即停；补出的合法续棋不唯一。

\ex{Monte Carlo 与 MCTS 简查（课件 05 p.73--74）}
在 $[-1,1]^2$ 均匀抽 N 点，K 点落单位圆内，则 $\pi\approx4K/N$，不是 K/N。MCTS 四步：选择（平衡探索/利用）→扩展→模拟 rollout→回传访问次数与累计收益。教学变式：访问 10 次、收益 6，新增收益 1 后变为 11 次、收益 7，均值 7/11；保持收益视角一致。

\clearpage\section{CSP、删域与局部修复}\label{sec:csp}
\src{课件 06 p.14--63；HW4 CSP1--3}
CSP 由变量、域、约束组成；解为\textbf{完整且满足全部约束}的赋值。
\par\noindent\begin{tabularx}{\textwidth}{lY}
\toprule 名称&执行规则\\\midrule
MRV&选当前合法域最小的未赋值变量（不是最小数字值）。\\
Degree&通常用来打破 MRV 并列，选与最多未赋值变量有约束者。\\
LCV&对已选变量，选删掉其他域中值最少的候选值。\\
Forward checking&赋值后删未赋值邻居中不相容值，任一空域则回溯。\\
Arc consistency&每个源域值在目标域有支持；没支持则删源域值，继续传播。\\
Min-conflicts&从完整但有冲突赋值出发，选冲突变量，把它改成违反最少约束的值。\\\bottomrule
\end{tabularx}\par
回溯返回上一层时要恢复该层造成的删域。

\ex{建模公式：皇后与密码算术}
皇后每列一个：$Q_i\in\{1,\ldots,n\}$；对 $i<j$，$Q_i\ne Q_j$，$|Q_i-Q_j|\ne|i-j|$。候选数 $n^n$；若预设行也互异才是 $n!$。解例 4 皇后 $(2,4,1,3)$。

TWO+TWO=FOUR：六字母 AllDifferent，数字 0--9，首位 T,F 非零，进位 $X_1,X_2,X_3\in\{0,1\}$：
\[
2O=R+10X_1,\quad X_1+2W=U+10X_2,\quad
X_2+2T=O+10X_3,\quad X_3=F=1.
\]
进位不属于字母 AllDifferent，不额外添加不同进位必须不同的条件。

\ex{MRV + degree + FC 全表（HW4 CSP1）}
边 $E=\{ab,ac,ad,bc,be,cd,ce,de\}$，端点不等。先做一元约束：$D_a=\{1,2\},D_b=\{1,2,3\}$，其余 $\{1,2,3,4\}$。规则 MRV→未赋值 degree→字母序；值升序。
\begin{center}\small\begin{tabular}{c|c|ccccc}\toprule
步&赋值&$D_a$&$D_b$&$D_c$&$D_d$&$D_e$\\\midrule
0&--&12&123&1234&1234&1234\\
1&a=1&--&23&234&234&1234\\
2&b=2&--&--&34&234&134\\
3&c=3&--&--&--&24&14\\
4&d=2&--&--&--&--&14\\
5&e=1&--&--&--&--&--\\\bottomrule\end{tabular}\end{center}
表中 24 表示集合 $\{2,4\}$；横线表示已赋值。答案 $\ans{(a,b,c,d,e)=(1,2,3,2,1)}$，无需回溯。d=4 也能完成，但不是本例数值升序选择。
\checkit{逐条检查八条边，不存在 b--d 边，所以 b=d=2 允许。a 与 e 不相邻，均为 1 允许。}

\ex{LCV、弧一致性与 AC-3（教学变式）}\label{ex:ac}
\textbf{LCV：}X 域 12，与 Y 域 1、Z 域 12 都要求不同。X=1 删 Y 的 1（空域）和 Z 的 1；X=2 只删 Z 的 2，故选 2。

\textbf{弧方向：}X$\to$Y 检查 $\forall x\in D_X,\exists y\in D_Y:C(x,y)$。若 X$\ne$Y，$D_X=\{R,B\},D_Y=\{B\}$，X 的 B 无支持，删 B 得 $D_X=\{R\}$。

\textbf{传播：}A$\ne$B、B$\ne$C，域 $D_A=\{R\},D_B=\{R,B\},D_C=\{B\}$。B$\to$A 删红，B 只剩蓝；C$\to$B 删蓝，C 空，当前分支无解。

AC-3：队列存有向弧；取 X$\to$Y 做 REVISE；若 $D_X$ 变化，把其他邻居的 Z$\to$X 重新排队；直到空域失败或队列空。常见复杂度 $O(ed^3)$。
\checkit{只在支持不存在时删值；弧一致不等于全局可解。两色三角形每条弧都有支持，整体仍无解。}

\ex{固定红点染色与四色下界（HW4 CSP2--3）}
CSP2：左 L，上左 A、下左 B、上右 C、下右 D；边 LA,LB,AB,AC,AD,BD,CD。L=红，故 A、B 必为绿蓝且不同；取 A=绿，B=蓝，D 被 A、B 排除只能红，C 被 A、D 排除只能蓝。
\[
\ans{(L,A,B,C,D)=(R,G,B,B,R)}.
\]
CSP3：修改地图的 SA,Q,NSW,V 形成 $K_4$，至少四色；一组可行值 WA=2,NT=3,SA=1,Q=2,NSW=3,V=4，所以 $\chi=4$。下界需要 K4 或等价排除理由；只画四色方案只能证明至多四。

\ex{min-conflicts 冲突计数（课件 06 p.54--58；教学变式）}
4 皇后每列一枚，初始 $(1,1,1,1)$。选第 2 列，候选行 1,2,3,4 与其他三枚皇后的冲突数分别 $3,2,1,0$，移到行 4 得 $(1,4,1,1)$。这是一次修复，尚非解；继续直到每对满足行/对角约束。整盘冲突按\textbf{皇后对}计一次，不能把各皇后冲突数相加后不除 2。

\ex{树结构与 Zebra 表示简查}
树状二元 CSP：叶到根做父对子的弧一致，再根到叶选择与父相容值；非空域时可解，典型 $O(nd^2)$。含环图可 cutset 固定后解树，或簇分解后在共享变量上匹配；簇大时组合数仍指数。

Zebra 位置变量取 1--5，同类 AllDifferent；同房用等号，邻房用 $|pos(A)-pos(B)|=1$，绿紧接象牙右侧为 $pos(G)=pos(I)+1$。挪威人在1、邻蓝故蓝在2；第一间排除红/蓝/绿/象牙后黄；黄配 Kit Kats，再排除牛奶/咖啡/茶/橙汁，故挪威人喝水。

\clearpage\section{命题逻辑、范式、SAT 与归结}\label{sec:logic}
\src{课件 07 p.8--87；HW5 A--D}
\begin{center}\begin{tabular}{cc|cccccc}\toprule
p&q&$\neg p$&$p\land q$&$p\lor q$&$p\oplus q$&$p\to q$&$p\leftrightarrow q$\\\midrule
F&F&T&F&F&F&T&T\\F&T&T&F&T&T&T&F\\T&F&F&F&T&T&F&F\\T&T&F&T&T&F&T&T\\\bottomrule
\end{tabular}\end{center}
优先级 $\neg>\land>\lor>\to>\leftrightarrow$；n 个不同变量真值表 $2^n$ 行。
\par\noindent\begin{tabularx}{\textwidth}{lY}
\toprule 词语/形式 & 对应意义\\\midrule
$p$ only if $q$；p 充分于 q & $p\to q$；q 是 p 的必要条件。\\
$p$ if $q$ & $q\to p$。$p$ unless $q$ 通常为 $\neg q\to p$。\\
逆/否/逆否 & $q\to p$；$\neg p\to\neg q$；$\neg q\to\neg p$。原式与逆否等价。\\
永真/矛盾/可满足 & 全行真 / 全行假 / 至少一行真；可满足不等于永真。\\
等价 $F\equiv G$ & 每行同真值；$F\leftrightarrow G$ 永真。\\
蕴含 $KB\models Q$ & KB 真时 Q 必真；等价于 $KB\land\neg Q$ 不可满足。\\\bottomrule
\end{tabularx}\par

\ex{等价律与推理规则}
\[
\begin{aligned}
p\to q&\equiv\neg p\lor q,&\neg(p\to q)&\equiv p\land\neg q,\\
p\leftrightarrow q&\equiv(\neg p\lor q)\land(\neg q\lor p),\\
\neg(p\land q)&\equiv\neg p\lor\neg q,&\neg(p\lor q)&\equiv\neg p\land\neg q,\\
p\lor(q\land r)&\equiv(p\lor q)\land(p\lor r),\\
p\land(q\lor r)&\equiv(p\land q)\lor(p\land r).
\end{aligned}
\]
幂等 $p\lor p=p$；吸收 $p\lor(p\land q)=p$；$p\lor\neg p=T,p\land\neg p=F$。Modus Ponens：$p,p\to q\vdash q$；Modus Tollens：$\neg q,p\to q\vdash\neg p$；$p\lor q,\neg p\vdash q$；$p\to q,q\to r\vdash p\to r$。\textbf{由 q 和 $p\to q$ 不能反推 p。}

\ex{NNF、CNF、DNF 完整转换（课件 07 p.60--70）}
文字 literal 为 p 或 $\neg p$；term 为文字 AND；clause 为文字 OR。
\begin{itemize}
\item NNF：只有 $\neg,\land,\lor$，否定只作用原子。
\item CNF：外 AND，内 OR，即子句合取；DNF：外 OR，内 AND，即项析取。
\item 流程：去 $\to,\leftrightarrow$ → De Morgan 推否定 → 去双否 → 向目标方向分配并化简。
\end{itemize}
\textbf{教学变式：}$F=\neg(p\to q)\lor r$。
\[
F=\neg(\neg p\lor q)\lor r=(p\land\neg q)\lor r
=(p\lor r)\land(\neg q\lor r).
\]
中间式为 NNF 和 DNF，最后为 CNF。

\textbf{课件 p.68：}
\[
\begin{aligned}
(p\to q)\land(q\to p)
&=(\neg p\lor q)\land(\neg q\lor p)\\
&=(\neg p\land\neg q)\lor(\neg p\land p)\lor(q\land\neg q)\lor(q\land p)\\
&=\ans{(\neg p\land\neg q)\lor(p\land q)}.
\end{aligned}
\]
\textbf{真值表转主范式：}真行写 AND 最小项（真写正、假写负），再 OR；假行写 OR 最大项（真写负、假写正），再 AND。例如假行 $(T,F,T)$ 对应 $\neg p\lor q\lor\neg r$。

\textbf{课件 p.69--70：}$F=(x\lor y)\to(\neg x\land z)$，真行 000,001,011，故
\[
\operatorname{PDNF}=(\neg x\land\neg y\land\neg z)\lor(\neg x\land\neg y\land z)\lor(\neg x\land y\land z).
\]
假行 010,100,101,110,111，各取最大项 AND：
\[
\begin{aligned}
\operatorname{PCNF}={}&(x\lor\neg y\lor z)\land(\neg x\lor y\lor z)\land(\neg x\lor y\lor\neg z)\\
&\land(\neg x\lor\neg y\lor z)\land(\neg x\lor\neg y\lor\neg z).
\end{aligned}
\]
简化 CNF 为 $\neg x\land(\neg y\lor z)$。主范式中每项含全部变量，简化范式不必。
\checkit{De Morgan 要同时换联结词并否定每个文字；不能只把 AND/OR 对调。CNF 分配成 DNF 可能指数膨胀。2-SAT 可线性求解，3-SAT 为 NP 完全；不是已经证明没有多项式算法。}

\ex{真值树：展开、关闭与反例}
检查蕴含，放入全部前提和\textbf{结论的否定}。
\begin{center}\begin{tabular}{l|l}\toprule
同一分支加入&分成两个分支\\\midrule
$A\land B$: A,B&$A\lor B$: A / B\\
$\neg(A\lor B)$: $\neg A,\neg B$&$\neg(A\land B)$: $\neg A$ / $\neg B$\\
$\neg(A\to B)$: $A,\neg B$&$A\to B$: $\neg A$ / B\\
$\neg\neg A$: A&$A\leftrightarrow B$: (A,B) / ($\neg A,\neg B$)\\\bottomrule
\end{tabular}\end{center}
某分支含 p 与 $\neg p$ 则关闭；所有分支全闭才不可满足。完全展开的开放分支给出模型/反例。

\textbf{教学变式：}$p\lor q,p\to r,q\to r\models r$：加 $\neg r$，两蕴含迫使 $\neg p,\neg q$；$p\lor q$ 两枝都关闭。

\textbf{课件 p.50--52：}$B\to C,D\lor C,B\leftrightarrow\neg D$ 不蕴含 $\neg(B\land C)$，反例 B=T,C=T,D=F。只要有一行前提全真而结论假，蕴含即不成立。

\ex{归结与不可满足（HW5 B）}\label{ex:resolution}
\[
\frac{l\lor A\qquad\neg l\lor B}{A\lor B}.
\]
一次消一对互补文字，保留原子句并加入归结式。一般是\textbf{推出}，不是把原式等价替换掉。$p,\neg p\vdash\square$ 空子句为假；空子句集无约束为真。
\begin{center}\begin{tabular}{c|l|l}\toprule
编号&子句&来源/消去文字\\\midrule
1&$p\lor q\lor\neg r\lor s$&原题\\
2&$\neg p\lor r\lor s$&原题\\
3&$\neg q\lor\neg r$&原题\\
4&$p\lor\neg s$&原题\\
5&$\neg p\lor\neg r$&原题\\
6&$r$&原题\\\midrule
7&$\neg q$&3,6；r\\
8&$\neg p$&5,6；r\\
9&$\neg s$&4,8；p\\
10&$p\lor q\lor s$&1,6；r\\
11&$p\lor s$&10,7；q\\
12&$s$&11,8；p\\
13&$\square$&12,9；s\\\bottomrule
\end{tabular}\end{center}
不可满足；原子句 2 不必用到。若只是归结出 $r\lor\neg r$，那是永真，不能宣布矛盾。

\textbf{课件 p.39 的短例：}$p\lor q,p\lor\neg q$ 推 p；$\neg p\lor r,\neg p\lor\neg r$ 推 $\neg p$；再得空子句。

\ex{反例与最大满足数（HW5 A、C）}
A 原件横线读为否定：$[(a\lor b)\land(\neg a\lor c)]\to\neg(b\lor c)$。反例 $\ans{a=T,b=F,c=T}$：前件真、后件假。标准归结正确结论是 $b\lor c$，不能为“证明”把原题横线无声删去。

C 五子句 $p\lor\neg q,\neg p\lor q,q\lor r,q\lor\neg r,\neg q\lor\neg r$，取 $\ans{p=q=T,r=F}$ 全部真，最大数 $\ans{5}$。上界是总子句数，给出达到上界的赋值即可。

\ex{同时成立与真值表（HW5 D）}
实数 $x=y,y=z,z=-x$ 推 $2x=0$，唯一解 $x=y=z=0$。
逻辑 $\Phi=(A\to B)\land(B\to C)\land(C\to\neg A)$：
\begin{center}\begin{tabular}{ccc|ccc|c}\toprule
A&B&C&$A\to B$&$B\to C$&$C\to\neg A$&$\Phi$\\\midrule
T&T&T&T&T&F&F\\T&T&F&T&F&T&F\\T&F&T&F&T&F&F\\T&F&F&F&T&T&F\\
F&T&T&T&T&T&T\\F&T&F&T&F&T&F\\F&F&T&T&T&T&T\\F&F&F&T&T&T&T\\\bottomrule
\end{tabular}\end{center}
三组满足赋值 FTT,FFT,FFF；$\Phi\equiv\neg A\land(\neg B\lor C)$，可满足而非永真，选“某些情况同时成立”。

\ex{Horn 与规则链（课件 07 p.76--87；教学变式）}
Horn 子句最多一个正文字。$\neg A\lor\neg B\lor C\equiv(A\land B)\to C$；事实 C 是单文字规则。事实 A,B；规则 $A\land B\to C,C\to D$：前向链逐步新增 C,D；后向查 D→需 C→需 A 和 B→事实支持。前件必须全满足才能触发；未证明某事实不自动等于它为假。专家系统：知识库、工作记忆、推理机、规则议程；匹配→选规则→执行→更新。

\clearpage\section{计数、二项、多项与平滑}\label{sec:prob}
\src{课件 08 p.3--19、47；HW6 P1}
\[
\binom nk=\frac{n!}{k!(n-k)!};\quad
P(A)=\frac{|A|}{|\Omega|}\ \text{仅适用于等可能基本结果}.
\]
\par\noindent\begin{tabularx}{\textwidth}{lY}
\toprule 模型 & 公式与条件\\\midrule
Bernoulli & $X\in\{0,1\},P(X=x)=p^x(1-p)^{1-x}$。\\
Binomial & n 次独立同概率试验，$P(K=k)=\binom nkp^k(1-p)^{n-k}$，$E[K]=np,Var(K)=np(1-p)$。\\
Multinomial & $\sum_i n_i=n,\sum_i p_i=1$，$P(\mathbf n)=\dfrac{n!}{\prod n_i!}\prod p_i^{n_i}$。\\
有序序列 & 每个位置类别已指定时直接 $\prod_i p_i^{n_i}$，不乘多项系数。\\
MLE & $\widehat p_i=n_i/N$；多批先合并计数和样本量。\\
Laplace add-one & $\widetilde p_i=(n_i+1)/(N+K)$，K 为该分布可能类别数。\\\bottomrule
\end{tabularx}\par
无放回抽样可能不独立；按题目抽样规则选模型，不自动套二项/多项。

\ex{硬币与合并估计（课件 08 p.6、11--12）}
10 次公平硬币恰好 5 正：$\binom{10}5/2^{10}=\ans{63/256}$；恰好 4 或 6 正：$2\binom{10}4/2^{10}=\ans{105/256}$。期望 5 不表示每次一定 5。

三批 4/10、4/8、7/12，合并 $\widehat p=(4+4+7)/(10+8+12)=\ans{1/2}$，不对比例简单平均。多项相同真实 p 且批次独立时，计数可加成 $\mathrm{Multinomial}(n+m,p)$；p 不同时不能直接套此分布性质。

\ex{多项组成与最可能组成（HW6 P1，p.1）}
已知学生五类比例 $p=(.24,.22,.20,.18,.16)$，四次独立抽取。原题要求 Laplace，但未给原样本总量 N。定义
\[
q_i=\frac{Np_i+1}{N+5};\qquad
P(\mathbf n)=\frac{4!}{\prod n_i!}\prod_i q_i^{n_i}.
\]
\begin{center}\begin{tabular}{l|l}\toprule
事件&计算结果\\\midrule
全四名大四&$\ans{q_4^4}$\\
两名大一、两名大二&$\ans{6q_1^2q_2^2}$\\
最可能组成&$\ans{(1,1,1,1,0)}$，概率 $\ans{24q_1q_2q_3q_4}$\\
最不可能组成&$\ans{(0,0,0,0,4)}$，概率 $\ans{q_5^4}$\\\bottomrule
\end{tabular}\end{center}
最大值理由：若某类 i 重复且 j 未用，移一个后概率乘 $n_iq_j/q_i$。本题 $q_i/q_j\le1.5$ 且 $n_i\ge2$，比值大于 1，故最大组成四类各一个，再取最大的四类。最小值：任一概率乘积至少 $q_5^4$，系数至少 1，全部选第五类达到下界。

\textbf{示范假设 N=100（非题给事实）：}$q=(25,23,21,19,17)/105$。四问对应
\[
\left(\frac{19}{105}\right)^4,\quad
6\left(\frac{25}{105}\right)^2\left(\frac{23}{105}\right)^2,\quad
\frac{24\cdot25\cdot23\cdot21\cdot19}{105^4},\quad
\left(\frac{17}{105}\right)^4.
\]
分别约 .00107215, .01632036, .04529964, .00068713。若不给 N，\textbf{以含 N 公式为正式结果}。不平滑时分别为 .00104976, .01672704, .0456192, .00065536。
\checkit{百分数并不自动给出 N=100；组成忽略次序时需要系数 6 或 24；最可能组成不一定全选最大 p 类。}

\ex{比较组成概率的手算比值（课件血型例）}
从 i 类移一个到 j 类，
\[
\frac{P(n_i-1,n_j+1)}{P(n_i,n_j)}=\frac{n_i}{n_j+1}\frac{p_j}{p_i}.
\]
血型 p=(.44,.42,.10,.04)，10 人组成 (6,2,1,1)：
\[
P=\frac{10!}{6!2!}(.44)^6(.42)^2(.10)(.04)\approx .01290.
\]
课件 08 p.19 说 (4,4,1,1) 概率 .0294 为最大；\textbf{若允许零计数，真正最大为 (5,4,1,0)，概率约 .0646593}；若附加每类至少一人，(4,4,1,1) 才是受限最大。计数总数小可枚举所有非负组成（此例 $\binom{13}3=286$），不可只比较一个候选就宣布全局最大。

\clearpage\section{Bayes 与 Naive Bayes 分类}\label{sec:bayes}
\src{课件 08 p.20--49；HW6 P2--P3}
\[
P(H|E)=\frac{P(E|H)P(H)}{P(E)},\quad
P(E)=\sum_jP(E|H_j)P(H_j).
\]
$H_j$ 互斥且穷尽。prior=$P(H)$；likelihood=$P(E|H)$；evidence=$P(E)$；posterior=$P(H|E)$。\textbf{先写条件方向，再代入。}

\ex{检测阳性与后验（课件 08 p.30--35）}
V=感染，$P(V)=.04$；灵敏度 $P(+|V)=.95$；特异度 $P(-|\neg V)=.99$，故假阳性 .01。
\begin{center}\begin{tabular}{c|c|c|c}\toprule
分支&先验&阳性条件概率&联合权重\\\midrule
V&.04&.95&.038\\$\neg V$&.96&.01&.0096\\\bottomrule\end{tabular}\end{center}
\[
P(V|+)=\frac{.038}{.038+.0096}=\ans{95/119}\approx.7983.
\]
不能把 $P(+|V)$ 的 .95 直接作为 $P(V|+)$。\checkit{归一化两类后验和为 1；注意特异度需先取补才能得到假阳性。}

\ex{答对后未准备（HW6 P2，p.2）}
U=未准备，R=答对。$P(U)=1/4,P(\neg U)=3/4$；$P(R|U)=1/4,P(R|\neg U)=9/10$。
\[
P(U|R)=\frac{(1/4)(1/4)}{(1/4)(1/4)+(3/4)(9/10)}
=\frac{1/16}{59/80}=\ans{5/59}.
\]

\ex{朴素贝叶斯的估计、分数与后验}
给定类别 c 后属性条件独立：
\[
P(\mathbf x|c)=\prod_jP(x_j|c),\quad s_c=P(c)\prod_jP(x_j|c),\quad
\widehat c=\arg\max_c s_c,\quad P(c|\mathbf x)=\frac{s_c}{\sum_{c'}s_{c'}}.
\]
MLE 条件概率 $n_{j,v,c}/N_c$；add-one：$(n_{j,v,c}+1)/(N_c+K_j)$，$K_j$ 是\textbf{该属性的取值数}。若先验也平滑，$P(c)=(N_c+1)/(N+\#classes)$。明确哪些分布做平滑，不随意混用。

\ex{动物分类全计数（HW6 P3，p.3）}
原表 20 例，M 哺乳=7，N 非哺乳=13。属性顺序 Birth,Fly,Water,Legs；\textbf{作业可见查询为四个 yes}，课件旧查询是 yes,no,yes,no。
\begin{center}\small\begin{tabular}{c|c|c|c|c}\toprule
类&Birth yes/no&Fly yes/no&Water yes/no/sometimes&Legs yes/no\\\midrule
M&6/1&1/6&2/5/0&5/2\\N&1/12&3/10&3/6/4&9/4\\\bottomrule\end{tabular}\end{center}
\[
\begin{aligned}
s_M&=\frac7{20}\frac67\frac17\frac27\frac57=\frac3{343},\\
s_N&=\frac{13}{20}\frac1{13}\frac3{13}\frac3{13}\frac9{13}=\frac{81}{43940}.
\end{aligned}
\]
$s_M>s_N$，判 M；归一化
\[
P(M|\mathbf x)=\frac{3/343}{3/343+81/43940}=\ans{43940/53201}\approx.825924.
\]
若全部 add-one，Water 有三值，其余二值：
\[
\widetilde s_M=\frac8{22}\frac79\frac29\frac3{10}\frac69,\qquad
\widetilde s_N=\frac{14}{22}\frac2{15}\frac4{15}\frac4{16}\frac{10}{15}.
\]
归一化 $P(M|x)=\ans{10/13}$，仍判 M。
\checkit{每个属性按其类别子表计数；分母不是全体 20；score 还需归一化才能叫后验。属性数 4 不是 Water 平滑分母加 4 的理由。}

\ex{银行分类与零概率（课件 08 p.43--47）}
标签是是否违约，No 类7例、Yes类3例。按原始表重数：
\begin{center}\small\begin{tabular}{c|c|c|c}\toprule
类&Home yes/no&Marital single/married/divorced&Income low/mid/high\\\midrule
No&3/4&2/4/1&3/1/3\\Yes&0/3&2/0/1&0/3/0\\\bottomrule\end{tabular}\end{center}
收入 low$<80$K、mid 80--101K、high$>101$K。p.45 把 married/divorced 数字调换，\textbf{以原表计数为准}。

教学查询 home=no,single,mid：
\[
s_{No}=\frac7{10}\frac47\frac27\frac17=\frac4{245},\quad
s_{Yes}=\frac3{10}\frac33\frac23\frac33=\frac15.
\]
判 Yes（违约），后验 $P(Yes|x)=49/53$。如查询 home=yes，未平滑 Yes 类整项为零；add-one 条件变 $(0+1)/(3+2)=1/5$，避免零概率。

\ex{连续属性与维数灾难简查（课件 08 p.23--29）}
连续 x 用密度 $f(x|c)$，不能把精确点概率当非零。正态
\[
f(x|c)=\frac1{\sigma_c\sqrt{2\pi}}e^{-(x-\mu_c)^2/(2\sigma_c^2)},\quad s_c=P(c)f(x|c).
\]
身高例 x=1.65，男/女均值 1.67/1.61、共同标准差 .05、先验 .7/.3，加权密度约 5.156/1.738；归一化男约 .748。无计算器可保留指数式。

每维 k 格，联合格子 $k^d$；在单位 d 维立方体覆盖体积 s 的小立方体边长 $s^{1/d}$。d=10，s=.1 时 .794，s=.01 时 .631（课件 p.28 对调）。Naive Bayes 通过条件独立减少需要估计的联合参数，但这是模型假设。

\clearpage\section{联合分布、独立性与 Markov 链}\label{sec:markov}
\src{课件 09 p.4--36、106、108；HW6 P4--P5}
\[
P(x)=\sum_yP(x,y),\quad P(x|y)=P(x,y)/P(y),\quad P(x,y)=P(x|y)P(y).
\]
独立：$P(x,y)=P(x)P(y)$ 对全部 x,y 成立。条件独立：$P(x,y|z)=P(x|z)P(y|z)$，等价于 $P(x|y,z)=P(x|z)$（所需条件概率有定义时）。条件独立与无条件独立互不自动推出。

\ex{联合表与 chi-square（课件 09 p.4--10）}
\begin{center}\begin{tabular}{l|rrr|r}\toprule
Game / Snack&Pizza&Chips&Cookies&合计\\\midrule
Poker&10&3&12&25\\Trivial Pursuit&8&14&7&29\\Monopoly&14&17&7&38\\Wii Bowling&12&7&4&23\\\midrule
合计&44&41&30&115\\\bottomrule\end{tabular}\end{center}
$P(Poker,Pizza)=10/115$；$P(Poker)=25/115$；$P(Pizza|Poker)=10/25$，反向 $P(Poker|Pizza)=10/44$。

$H_0$ 为行列变量独立，期望\textbf{频数} $E_{ij}=R_iC_j/N$；
\[
\chi^2=\sum_{i,j}\frac{(O_{ij}-E_{ij})^2}{E_{ij}},\quad df=(r-1)(c-1).
\]
$E_{11}=25\cdot44/115=220/23$；全表 $\chi^2\approx14.508$，df=6，课件给5\%临界值12.592，拒绝独立。若另一情形统计量10.281，则不拒绝，\textbf{不等于证明独立}。抽样波动可令经验频率不精确分解。

\textbf{教学变式：}$O=\left[\begin{smallmatrix}8&2\\2&8\end{smallmatrix}\right]$，行列总计各10，N=20，四格 E=5；$\chi^2=4\cdot9/5=36/5>3.84$（df1），拒绝独立。Pearson 近似需独立观测和适当期望频数。

\ex{Markov 假设、方向和序列概率}
\[
P(X_{t+1}|X_t,\ldots,X_1)=P(X_{t+1}|X_t),\qquad a_{ij}=P(X_{t+1}=j|X_t=i).
\]
行 i 当前、列 j 下一状态；每行非负且和1；齐次指 A 不随时间改变，不等于边缘分布已稳态。还需初始分布 $\mu$。
\[
P(x_1,\ldots,x_T)=\mu(x_1)\prod_{t=1}^{T-1}a_{x_t,x_{t+1}},\quad
p_{t+1}=p_tA,\quad p_{t+k}=p_tA^k.
\]
使用\textbf{行向量}；$p_1=\mu$ 时 $p_t=\mu A^{t-1}$。多步概率对中间状态求和，特定序列按路径相乘。

\ex{17 天天气估计转移矩阵（HW6 P4）}\label{ex:transition}
顺序(S,C,R)，观察 $S,S,R,C,C,R,R,S,S,S,R,C,C,R,S,S,C$。
\begin{center}\begin{tabular}{c|rrr|r}\toprule
从/到&S&C&R&出发总数\\\midrule
S&4&1&2&7\\C&0&2&2&4\\R&2&2&1&5\\\midrule
合计&6&5&5&16\\\bottomrule\end{tabular}\end{center}
\[
\widehat A_{ij}=\frac{N_{ij}}{\sum_jN_{ij}},\qquad
\ans{\widehat A=\begin{pmatrix}4/7&1/7&2/7\\0&1/2&1/2\\2/5&2/5&1/5\end{pmatrix}}.
\]
\checkit{17天16个转移；末日 C 无后继，不作为出发计数。课件 p.18 的2/5与可见序列不符，C行分母应4。无出发观测的行无法直接估计，不可填零行冒充概率分布。}

\ex{一步、多步传播与稳态（课件 09 p.25--31）}
两状态(S,R)，$A=\left[\begin{smallmatrix}.9&.1\\.3&.7\end{smallmatrix}\right]$，已知第一天 S。
\begin{center}\begin{tabular}{c|cc}\toprule
日&p(S)&p(R)\\\midrule
1&1&0\\2&9/10&1/10\\3&21/25&4/25\\4&201/250&49/250\\\bottomrule\end{tabular}\end{center}
第三天晴 $.9^2+.1\cdot.3=.84$，不只包含连晴的.81。
稳态 $\pi=\pi A,\pi_S+\pi_R=1$，$.1\pi_S=.3\pi_R$，得 $\ans{\pi=(3/4,1/4)}$。

有限不可约非周期链保证任意初始分布收敛到唯一稳态；存在稳态不等于一定收敛，周期链可能振荡。

\ex{三状态矩阵的统一演算（HW6 P5；课件 p.32--36）}
统一顺序(S,C,R)，
\[
A=\begin{pmatrix}4/5&3/20&1/20\\19/50&3/5&1/50\\3/4&1/20&1/5\end{pmatrix}.
\]
\textbf{注意：}p.32 把第二/第三状态称 rainy/snowy，后续作业称 cloudy/rainy；先确认行列命名，不与上一题的估计矩阵混用。

教学初始分布 $p=(1/2,1/3,1/6)$：
\[
\begin{aligned}
p'(S)&=\tfrac12\tfrac45+\tfrac13\tfrac{19}{50}+\tfrac16\tfrac34=391/600,\\
p'(C)&=\tfrac12\tfrac3{20}+\tfrac13\tfrac35+\tfrac16\tfrac1{20}=17/60,\\
p'(R)&=\tfrac12\tfrac1{20}+\tfrac13\tfrac1{50}+\tfrac16\tfrac15=13/200.
\end{aligned}
\]
检查 $(391+170+39)/600=1$。

稳态联立归一化与两个独立方程：
\[
20\pi_S=38\pi_C+75\pi_R,\quad40\pi_C=15\pi_S+5\pi_R.
\]
第二式给 $\pi_C=(3\pi_S+\pi_R)/8$，代入得 $46\pi_S=638\pi_R$，所以
\[
\ans{\pi=(638,245,46)/929}.
\]
\checkit{$\pi A=\pi$，且分子和929。}

\ex{三天晴天与四天最可能序列}
\textbf{HW6 P5：}固定从 t 开始连续三天晴：
\[
\ans{P(SSS)=p_t(S)(4/5)^2}.
\]
已知首日晴则条件概率16/25；若明确稳态，则 $\ans{10208/23225}$。原题“可任意一天开始”\textbf{不自动给初始分布}，先给一般式。固定三天窗口不同于一周内任意处至少出现一次三连晴，后者窗口重叠不可直接乘。

\textbf{课件 p.34：}指定稳态开始，最可能四天 SSSS，概率 $\ans{(638/929)(4/5)^3=40832/116125}$。本例最大初始项为 S、最大边为 SS=.8，SSSS 同时达到各因子上界。

一般情况：$\delta_1(j)=\mu_j$，$\delta_{t+1}(j)=\max_i\delta_t(i)a_{ij}$，记最优前驱后回溯。\textbf{传播用 sum，最可能序列用 max}；不能拼接每天边缘最大状态就当最可能序列。

\ex{连续段与首次晴天等待（课件 p.35--36）}
已知今天晴，$q=a_{SS}<1$，额外连续晴天数K从0开始，$P(K=k)=q^k(1-q)$：
\[
E[K]=q/(1-q),\qquad E[L]=1/(1-q)\quad(L=K+1,\text{含今天}).
\]
q=4/5时，额外4天，含今天5天；若额外条件“明天仍晴”，该条件下的额外段长才为 $1/(1-q)$。不能直接照抄 p.35 的“more”。

首次到 S 的等待 $h_S=0$，非S状态 $h_i=1+\sum_ja_{ij}h_j$（可达且期望有限）。本例
\[
h_C=1+\tfrac35h_C+\tfrac1{50}h_R,\quad h_R=1+\tfrac1{20}h_C+\tfrac15h_R.
\]
化为 $20h_C-h_R=50,-h_C+16h_R=20$，解
\[
\ans{h_C=820/319,\quad h_R=450/319}\ \text{天}.
\]
从第三状态 R 出发，明天晴3/4，后天晴
\[
(A^2)_{RS}=\tfrac34\tfrac45+\tfrac1{20}\tfrac{19}{50}+\tfrac15\tfrac34=\ans{769/1000}.
\]
若要求\textbf{首次}后天晴，排除经 S 的中间路线，仅算 $a_{RC}a_{CS}+a_{RR}a_{RS}$。

\clearpage\section{作业题号索引与易错条件}\label{sec:crossref}
\begin{longtable}{p{.18\textwidth}p{.57\textwidth}r}
\toprule 原题 & 方法和查询入口 & 本册页\\\midrule\endhead
HW1 INTRO1--3 & 理性、人的任务局限、XOR & \pageref{sec:intro}\\
HW2 BFS1 & 八数码逐层计数 & \pageref{sec:search}\\
HW2 DFS1--4 & 栈、时间戳、森林、回溯口径 & \pageref{sec:search}\\
HW2 DFS5 & 逆退出时间与拓扑序检查 & \pageref{ex:topo}\\
HW2 Dijkstra1 & 非负边松弛，原题 bh=1 & \pageref{ex:dijkstra}\\
HW2 AStar1--2 & OPEN 更新，h(B)=16/6 & \pageref{sec:astar}\\
HW2 AStar3 & 连续状态与可见图 & \pageref{sec:astar}\\
HW3 Game1--5 & 前瞻计数、终局、回传、开放线 & \pageref{sec:games}\\
课件 05 p.66 & 指定顺序的两棵 alpha--beta 树 & \pageref{ex:ab}\\
HW4 CSP1--3 & MRV/degree/FC 与染色 & \pageref{sec:csp}\\
课件 06 p.40--58 & 弧一致性与 min-conflicts & \pageref{ex:ac}\\
HW5 A--D & 反例、归结、最大满足、真值表 & \pageref{sec:logic}\\
课件 07 p.60--82 & 范式、真值树、Horn 与链推理 & \pageref{sec:logic}\\
HW6 P1 & 多项组成、未知 N 的 Laplace & \pageref{sec:prob}\\
HW6 P2 & 答对后未准备，Bayes & \pageref{sec:bayes}\\
HW6 P3 & 可见查询四个 yes，按原表计数 & \pageref{sec:bayes}\\
HW6 P4--P5 & 转移估计、三天窗口与稳态 & \pageref{ex:transition}\\
课件 09 p.4--36 & 联合、卡方、传播、序列、等待 & \pageref{sec:markov}\\\bottomrule
\end{longtable}
\subsection*{开算前的十个条件}
\begin{enumerate}
\item 搜索根深度、生成/检查/扩展口径、去重、邻接与平局次序。
\item BFS 等步代价；Dijkstra 非负边；目标按规定时机停止。
\item DFS 发现序不是路径；进入退出均加1；有向图可有交叉边；拓扑序先排除环。
\item A* 松弛看 g，选点看 f，parent 同步更新；可采纳与一致分开。
\item 博弈看谁即将行动，终局停树；剪掉子树不虚填精确值。
\item MRV 选变量，LCV 选值；回溯恢复域；弧 X$\to$Y 删 X 无支持值。
\item 蕴含真前件假后件时假；反例前提全真；归结一次消一对；CNF/DNF外层别反。
\item 指定次序不乘组合系数，无序组成要乘；平滑需要样本量和取值数。
\item Bayes 先写条件方向；Naive Bayes 乘类先验；score 再归一化才是后验。
\item Markov 写状态顺序和初始分布；T天T$-$1次转移；任意日起点不自动稳态。
\end{enumerate}
\subsection*{来源与使用边界}
编号为复习顺序，不是官方 Lecture 编号。原课件在课程 slides/编号复习版；教师作业原件在 assignments/hw01--hw06/requirements。按复习计划与考试通知取范围，完整 Markdown 讲义另解释课件扩展部分。旧通知不改变本次考试安排。

本册按原件可见内容重算，明确修正了部分累计代价、概率、频数和范式问题；详细辨析见对应讲义。个人可在页边追加实际错题。Bayesian Network、HW6 P6--P7 未收入。

\end{document}
